\chapter{Introduction}

\section{Motivation}

Credit scoring is a consequential application of statistical and machine-learning
models in the financial sector: scoring outcomes can influence access to loans and
the conditions under which credit is offered, thus affecting an important dimension
of economic participation
\citep{bahloolPerformanceFairnessExplainability2026}. Model choice in this context
involves a tension. Flexible boosted-tree ensembles can represent nonlinear
relationships and interactions, but their fitted decision logic is not directly
inspectable and is therefore commonly explained post hoc. Inherently interpretable
models expose their decision structure more directly, but are often assumed to
sacrifice predictive performance
\citep{rudinStopExplainingBlack2019a, dessainCostExplainabilityAI2023}.

A recent systematic literature review questions how firmly this presumed trade-off is
established. It finds that predictive performance and explainability in AI-based credit
scoring are frequently assessed separately and that the
performance--explainability trade-off is more often assumed than directly quantified
\citep{bahloolPerformanceFairnessExplainability2026}. Existing direct comparisons
indicate that differences between interpretable and black-box models can be small and
depend on the dataset, preprocessing and model specification. Greater complexity
therefore does not necessarily produce a substantial predictive advantage
\citep{dessainCostExplainabilityAI2023,
amekoeExploringAccuracyInterpretability2024b}.

This issue also has regulatory relevance in the European Union. Under
Article~6(2) and Annex~III(5)(b) of Regulation~(EU)~2024/1689, creditworthiness
assessment is listed among the high-risk use cases for systems that fall within the
Act's definition of AI. The Act prescribes no model family but establishes requirements
concerning risk management, documentation, transparency, accuracy and human oversight
for covered high-risk systems. The GDPR applies independently where personal data are
processed and provides additional safeguards concerning certain automated individual
decisions \citep{mullerAIbasedCreditScoring2025, EUAIAct2024}.

To examine whether boosted-tree ensembles provide a consistent predictive or
profit-based advantage, this thesis compares two intrinsically interpretable models,
logistic regression (LR) and the Explainable Boosting Machine (EBM), with the
post-hoc-explained XGBoost and CatBoost ensembles. The EBM provides the central
glass-box case because it retains an additive, directly inspectable structure while
representing nonlinear effects and selected pairwise interactions
\citep{louAccurateIntelligibleModels2013,
noriInterpretMLUnifiedFramework2019}. Using two boosted-tree implementations also
examines whether recorded patterns recur within that model family.

Predictive performance and Expected Maximum Profit (EMP)
\citep{verbrakenDevelopmentApplicationConsumer2014} form the primary empirical
dimensions. Attribution-ranking stability and cross-model agreement provide supplementary
evidence on the resulting global explanation outputs, covering the reproducibility
and convergence of those outputs.

\section{Research Question and Objectives}\label{sec:rq}

The central research question is: \textit{Under the empirical conditions studied,
how do intrinsically interpretable credit-scoring models, particularly the EBM,
compare with post-hoc-explained boosted-tree ensembles in predictive and profit-based
performance, and what do the observed differences imply for model choice in the EU
regulatory context?}

The question is addressed in three empirical steps. First, the four pipelines are
compared across four public credit datasets in discrimination, classification performance and EMP. Second, the case-sampling stability
and cross-model agreement of their global source-feature attribution rankings are
examined. Third, training-size sensitivity is evaluated on Home Credit and Lending
Club, supplemented by out-of-time validation concerning Lending Club. The regulatory component situates these outcomes within the requirements that apply
to credit scoring under EU law, establishing what the empirical evidence can
contribute to a system-level assessment that reaches beyond model architecture. 

Chapter~\ref{ch:background} develops the conceptual, methodological and regulatory
background and derives the research gap from prior work, and
Chapter~\ref{ch:method} specifies the datasets, the four pipelines, the evaluation
measures and the prespecified analysis plan. Chapter~\ref{ch:results} reports the
performance, attribution and sensitivity results, Chapter~\ref{ch:discussion}
interprets them and states the limitations of the design, and
Chapter~\ref{ch:conclusion} answers the research question.
